\section{孙吴的地方社会与江河世界}
\label{reader:wu-river-society:start}

一枚写有户名或米数的木简，并不会自动把三国江南变成一幅井然有序的税收图。
它首先意味着一间县廷里有人要辨识这个名字，把田、租、役或转递的责任接到
另一根简上，再交给能够收、藏、发、核的人。出土于长沙走马楼的孙吴简牍，正
是在这个尺度上把临湘县的行政劳动留给后人。与《三国志》里帝号、北伐和宫廷
政变相比，这些木片显得迟缓；但船粮、道路工程和守江的命令若不能变成这种迟
缓的手续，便没有办法抵达田野与渡口。

这不是要用临湘替孙吴画出一张全境统计表。简牍的保存由建筑、火灾、埋藏和发
掘的偶然共同决定；留下来的是某些机构曾经需要处理的事务，不是全吴每一县、
每一年都照此运作的副本。它们没有给出一个可以安全加总的人口数，也没有使我
们知道建业、会稽、武昌或交州各以何种税率征收。可是，正因为材料只照亮一个
县廷的近景，它迫使读者不把“江东”当作一块由长江自然供养的腹地：国家能力
必须逐户、逐田、逐仓门地被做出来，也会在这些接口上失去效力。

\subsection{文书不是从一座完整档案库里取出的}

走马楼的发现首先是一项考古事实，不是孙吴主动留给后世的一部县志。被掩埋
的简牍在发掘中以残片、散片和相互关系不总是完整可复原的状态重见天日；今人
所读的门类、文书链和机构名称，都要经过释读、缀合与比较。因而，“一根简”
往往不是一份自足的报告。它可能只保存了一个名单的尾部、一项收发的提示、某
次核验的数目，或者一组原本要和别简一道使用的标记。将它们编排为可读材料，
是必要的工作，却不应抹去它们原来并不为讲述地方社会而写的事实。

这个发现过程改变了史料在叙事中的位置。《吴主传》把一年压缩成皇帝的诏令、
出兵、任免和祭祀，读者可以沿年号较快地前行；木简则拒绝这种速度。它们不必
注明与宫廷纪年相同的时刻，所关心的是某一笔租、某一户、某种役何时需要补入
程序。两类材料不能互相校正成一张无缝年表：帝纪的“省徭役”不必在简中留下
可辨的回声，简中的一项勾核也未必能归到某次全国政策。把它们放在同一卷中，
不是为了让地方材料替朝廷叙事增加几个例证，而是为了让读者同时感到国家有两
种时间：诏书的时间，以及一件事务迟迟不能结案的时间。

考古材料还限制了我们的视野。能被保存并被辨读的，多半是与管理有关的文字；
不会写字、没有机会进入县廷的人，仍然大多不以自己的声音出现。也不能因为一
些田、租、役的术语保存下来，便假设临湘人人都以同样方式理解官府。反过来说，
正史中没有姓名的船工、佃作者、仓吏和被征者，也不能被当成不存在。材料的偏
向不是一个可以用想象填平的缺口，而是阅读三国地方社会时应一直带在身边的限
制：我们看见的是国家碰到他们时留下的边缘。

\subsection{临湘县廷：把人、田与责任暂时拴在一起}

走马楼简中可辨的户、田、租米、吏役和勾核，并不是几个可以直接翻成现代财
政栏目词。它们原来属于不同时间的收发、登记和复核。书手写下一个户名，
不等于此户从此不迁徙；保管者记录一批米，也不等于米已经无损地进了都城的仓。
恰恰相反，文书必须反复写、改、对照，才说明人、地和物不会安静地停在官府给
它们的位置上。疾病、逃亡、依附、土地变更、欠纳，或者只是一段水路难行，都
会让前一次分类在下一次差役到来时不再够用。

这里能够看见的家庭，不是带着完整生平走出史料的“普通人”。简牍很少让我们
知道某一户如何讨论婚姻、分家或离乡，也没有许可我们为一个名字补写悲欢。它
让家庭出现的方式更冷静：作为须被辨认、被追问、被分摊某项责任的单位。正是
这个不完整的出现很重要。孙吴朝廷可以在本纪中宣布征发、赐予或减省，县廷所
面对的却是哪些人可被列名，哪些田可与租相连，哪一项差役需要由谁补上。诏令
越过长江以后，先变成的并不是抽象的“社会控制”，而是让基层中介和官吏能够
继续工作的细目。

临湘的木简还改变了“户籍”一词的读法。户籍当然给国家提供了一个可想象的资
源库：它使征粮、派役、调人有了可追索的对象。但它也把风险压回家庭。被登记
的人不只是在册的人；他或她还可能因亲属、佃作关系、迁居或前一项未清的责任，
被卷入新的查验。现存简牍不足以告诉我们这些关系在全吴如何运作，更不能把一
县的格式宣称为孙权或孙皓的统一制度。它所容许的较小而坚实的判断是：江防的
粮与役并不从一个没有摩擦的“民间”流入军府，二者之间有不断登记、保管和争
议的行政现场。

\begin{sourcenote}[走马楼吴简看见的是一座县廷，不是一部孙吴总账]
木简使临湘县的户、田、租与吏役可见，因而能校正只凭帝纪理解国家的视角；它
们却不能提供全吴人口、亩数、税率或制度执行的一致图景。把简牍同《吴书》的
年序并读，能问的是朝廷资源为何需要经由地方手续；不能据此断言每个江东家庭
都承担同样的租役。
\end{sourcenote}

\subsection{田租不是一条从田里直通建业的线}

田地在文书中能够被识别，租米才可能被催收、储藏或转输。这条次序看似朴素，
却解释了为什么战争国家不能只靠一座都城的命令生活。稻米离开田面以后，还要
经过量、纳、仓、舟和再分配；其中任何一处延误，都可能使县廷面对的不是一项
完成的收入，而是一条需要继续追问的责任。我们没有连续账簿来画出孙吴粮食由
临湘抵达何处的全部路线，也不该把“租米”想成必然供给某一支水军。简牍的意
义并不在于补足这条不存在的总表，而在于让收取本身显出劳动和时间。

孙权在二二六年曾以“受田”和亲耕的语言表明与民共劳的姿态。这是《吴主传》
保留下来的统治者话语，能够说明朝廷愿意怎样表述田与劳作，不能证明皇帝亲耕
改变了各地的占有关系，亦不能替临湘的具体登记提供动机。将这段话与县廷简牍
并置，反而能看出两种尺度的距离：宫廷需要把农事说成德政与秩序，地方机构则
须在一片片可辨的田地和一笔笔可处理的责任中使粮食出现。二者并不互相否定，
但也绝非同一件事。

二五一年，孙吴朝廷下诏省徭役、减征赋、除民所患苦。它在年序中紧接着此前宫
廷冲突与军政紧张的时期，至少说明减轻负担已经成为能够被公开宣布的政治措辞。
不过，本纪的一句话没有留下哪个州郡何时收到诏书、哪些项目被停、欠负如何处
理，或家庭是否因而真正得到喘息。若把诏令直接当作社会结果，便会抹去走马楼
木简已经呈现的那些接口。比较稳妥的读法是把它看作一次承认：征发和负担会伤
及国家赖以取得人、粮与服从的地方关系；至于承认能否转成普遍的减免，材料没
有给出同样清楚的回答。

赤乌元年铸行大钱又从另一个方向显示压力。高额面值的货币安排需要金属、铸工、
官方许可以及愿意接受它的交易场景；朝廷的记载可确认这项尝试，却没有留下可
靠的兑换关系、物价系列或各地流通比例。它不能独自证明孙吴“富裕”或“通货
膨胀”，更不能替田租和役使算出货币化程度。大钱之所以应放回地方社会，是因
为国家在动员水军、工程和都城供给时，必须同时处理粮、劳力和可支付的媒介；
每一项安排都可能把不确定性传给下一个收取、交换或服役的人。

\subsection{东安郡与“山越”：山地不是一块等待填色的后方}

二二六年，《吴主传》记孙吴从三郡的“恶地”十县分置东安郡，以全琮为太守，
并称其“平讨山越”。短短一句把行政重划、山地行动和朝廷的胜利语言压在一
起。对建业而言，置郡意味着把一片难以稳固接入州郡秩序的空间重新命名、重新
配置官员；对沿山与水网生活的人而言，它也可能意味着军队、征发、道路和新的
地方中介进入旧有往来。可是“山越”不是一份能让我们读出共同政体、统一边界
或共同意志的自述材料，而是吴廷叙事使用的分类。

因此，不能从“平讨”二字推定所有相关山地群体立刻服属，也不能把他们写成一
支始终与吴对抗的民族军队。本纪没有保留他们内部如何决定迁徙、结盟、交涉或
避战的连续记录；山地地名、道路和社区的变化也不能靠一条纪事复原。能够确认
的是，孙吴的行政中心把这片区域视为需要军政处理的难题，而难题本身与土地、
人口和交通有关：若山地与沿江之间的人可以流动，县廷的户与役就不可能只是静
态名册；若水道和山口不易通行，置郡与派守又不会只是地图上的一笔。

全琮被放在这个位置，提醒人们孙吴的扩张不是只靠在江面赢得一次战役。地方官
须在军事行动、安置、征收和与既有地方势力的往来之间取得某种可持续的秩序。
史书把成败归入朝廷，未留下每座聚落承担的代价；我们无从断言哪些人被编入户
籍、哪些人留下或被迫迁走，更无从为“开发”写出一个单线的结果。但东安郡的
设置至少使江南的另一面显形：江河国家的边缘并不只在北岸，山地和内陆水网同
样会决定可征之人、可运之粮以及官府命令的尽头。

\begin{debate}[“山越”能告诉我们什么，不能替谁发言]
《吴主传》的称谓与结果来自孙吴的政治视角，可以确立二二六年置东安郡、派全
琮及相关军事行动的年序位置；它不能据此重建被归入“山越”的人群构成、语言、
疆界或长期政治选择。区域考古和历史地理能帮助提出山地、水网与聚落的问题，
却不能把沉默的人群化成一张已知的行政地图。
\end{debate}

\subsection{船旁的仓门：道路、邸阁与沿江劳作}

江防的船并不是在战事来临时才出现。船材要从山地和市场进入造修地点，铁器、
索具和桨手要被取得，粮食则须有仓门、码头和能够继续转运的人。港口不是一个
地图上永远亮着的点，而是堆卸、等候水位、辨认来船、保管货物与接收命令的多
种劳动暂时相会的地方。考古城址和河道研究可以让这些节点的长期空间条件较为
可想，却不能替正史编出一支舟师逐日经过的船表。江面优势从来不是河流赠予吴
廷的自然属性；它是许多难以写入本纪的维修、保管和信息传递累积出来的能力。

二四五年，陈勋率屯田与作士开凿句容中道，自小其通往云阳西城，又通会市、作
邸阁。这里的“通”尤其值得停下来。它不是抽象地“发展交通”：一条中道若能
连接水陆，便会改变粮、木材、军人和消息在何处换手；邸阁若确实建成，便为贮
放与转发提供了新的门槛。记载还列出三万人，这个数字作为《吴主传》的官修叙
事应被保留为动员规模的说法，不能被换成已核实的劳动力普查，也不足以重建每
一段工程的长度、死亡或地方摊派。

但正因数字不能被滥用，工程的选择才更加清楚。孙吴没有只等待战船在长江上解
决一切；它试图把屯田、作士、道路、市与邸阁拼成一条可供给都城和军队的通道。
这种做法暂时解决的是水陆交接的瓶颈，成本却不只在官府账上。被调出的屯田人
和作士不再完全按原来的节奏耕作或营造，沿线聚落要面对经过的人员、材料和新
的征调，仓储与市场又可能吸走附近劳力。哪些地方实际获益、哪些家庭被迫承担，
现有材料不能逐项回答；把这种不能回答保留在工程旁，才不会把“通会市”误读
为所有人共享的繁荣。

东兴、濡须和西陵等江防战事也应从这条日常供应线来理解。堤、渡、营垒和沿岸
据点能否发挥作用，取决于守军是否先得到消息、舟楫是否可用、粮食能否到岸。
它们并非使河流自动替吴军作战。对方也会利用同一片水面、风向和季节；二二四
年曹丕到广陵却因冰冻、舟师与渡江条件不具而撤回，正显示长江不是永恒不变的
屏障。水文条件会约束进攻者，也会约束守备者。沿江社会为军事提供的不是一个
现成后方，而是一条随季节、劳力与消息而时断时续的链条。

\subsection{濡须的船与一条不能只按将领命名的江面}

《吴主传》在二二二、二二三年的江淮战事中保存了几段很短的水上记载。曹休遣
臧霸以轻船进攻徐陵，吴方则令吕范等督数路舟军；翌年曹仁所部乘油船晨渡濡须
中洲。它们当然首先是战争叙事，所给出的兵数、战果与时间都来自各自政权的纪
录，不能让我们复原每艘船的载量、每个渡口的潮位，也不足以证明沿岸居民如何
看待这次来去。可是“轻船”“油船”“晨渡”这些词并非纯粹的修辞。它们提示
进攻与防守都要利用不同船型、能见度、岸线和中洲；江面并不是两支抽象军队相
遇的平板。

濡须的难处还在于消息。守方若在船队已到岸才得知来袭，即使岸上有兵也可能来
不及把船、粮与人调到同一个入口；若过早集中，又会使其他渡口、仓所和耕作失
去劳力。烽火、巡逻、商旅、地方官的急报和熟悉浅滩的人，未必都在正史中留下
姓名，但没有这些反复的传递，诸将的“督军”也只是一张不能移动的名册。这里
不能据沉默杜撰一套完整的预警制度，只能说战场叙事本身已显示出水上行动依赖
时间差，而时间差需要沿岸社会不断付出维护成本。

把这种成本放进地方社会，并非要把每个渔人或船工都写成军事人员。许多船只可
能服务于市场、渡运、收租或家庭迁徙，何时被国家征用、以何种条件补偿，现存
材料并不整齐。正是这种不确定性说明“水军国家”不是一个单独的军种名称。它
会在战时把船材、劳力、港口与水上知识拉进军令，也会在平时把同样的资源留在
田租、贸易和亲属生活之中。孙吴能够守住一些江面，并不表示它能无代价地占有
一切舟楫；守备所依赖的，始终是与地方用途竞争的有限物件和有限的人。

\subsection{句容中道怎样把市场接进防务}

陈勋开凿中道的记载之所以不能只当作工程新闻，还因为它把“会市”写进了同一
次行动。市场不是都城命令的终点，而是许多物品重新估价、换手和被人带走的地
方。道路接到云阳西城，邸阁又接到道路，意味着军府想在水陆之间建立一组较可
预测的转运节点；可是市中的交换未必完全服从军府的需要，仓储也会因水位、缺
粮、损耗、看守和临时征用而改变原先的用途。帝纪没有为我们留下每一间邸阁的
收支，不能将“作邸阁”直接译成一套已成功运行的物流体系。

这条工程还让人看见地方精英与基层吏员的作用，却不许我们轻易替他们安排立场。
要让道路、市与仓相接，官府需要能够认路、验收、征集材料、处理争执的人；这
些人可能来自既有的乡里关系，也可能因工程而重新取得地位。文献没有留下足够
材料，让我们把他们统称为支持或反对孙吴的一个集团。较有把握的是，工程把权
力落实为许多小决定：哪里可通，何物可藏，谁的劳力被调走，哪一次延误需要由
谁解释。地方社会不是在这些决定之外承受“国家影响”，而是国家每一次转运得
以继续的条件。

\subsection{从会稽望向外海：熟悉的航路并不等于海上统治}

二三〇年卫温、诸葛直奉命浮海求夷洲、亶洲，常被后人处理成一个急于标定现代
地名的问题。《吴主传》本身却先留下了一层更近的海上经验：会稽的货布能够到
来，东县人有海行之事，也有人遭风漂流到传称的亶洲。这样的句子不证明吴廷已
经掌握一条稳定的远洋航线；反而说明海边的货物、船人、传闻和意外漂移先于、
并且未必服从于朝廷的远航计划。港口中介和船工在这段材料中没有完整姓名和账
目，却不能因此被压缩成帝王远略的背景。

卫温、诸葛直的舟师离开近岸以后，要同风候、淡水、疾病、船况和补给竞争。同
一部本纪记船队未能如愿抵达，又得夷洲数千人而还，随后卫温、诸葛直因没有达
到预期被处置。它能确认朝廷确曾组织这次行动及其失败性的后果，不能让我们确
定夷洲、亶洲对应今日何地，不能核实传述中的航程与病死数，也不能把被带回者
的身份与遭遇写成已知。海上沉默不等候后人用国界填满。较可靠的结论是：吴廷
能把江面的人船集中为一次探索，却不能由此把外海变成可以持续征调和管理的水
域。

三年之后，孙吴又以张弥、许晏等携金宝珍货、九锡备物乘海赴辽东，回应公孙渊
的联络。这里的海路不是寻访岛屿，而是承担一项高风险外交：礼物、使者、随行
兵员和船只必须从江东港口离岸，越过北方海域，才能使远方的册拜成为有效行动。
后来公孙渊杀害吴使并转向魏，恰好显示海上到达不等于政治关系已经稳固。我们
不能根据一段朝廷叙事复原使团的准确泊地、人数或航程，更不应让“万兵”一类
记述替代实证；但这次反转让人看见，港口储备的珍货与舟船并不能消除对岸政权
重新计算利益的能力。

\begin{sourcenote}[海上消息有多种来处，不能合成一条精确航线]
《吴主传》同时保存了会稽海行、卫温远航和辽东使团等不同性质的记载。它们足
以显示孙吴与近海、岛屿传闻和东北外交均有接触，却不能证明一条连续、安全且
由吴廷控制的航道。夷洲、亶洲的地望、船队日程、人员损失和岛屿社会的内部情
况，都超出这些材料能够确认的范围。
\end{sourcenote}

\subsection{珠崖、儋耳与末期征调：远方也有自己的地方社会}

二四二年，聂友、陆凯受命率军讨珠崖、儋耳。这一纪事把孙吴的视线推到海南岛
及其周边的海峡，但它不能把那片空间写成从建业望去的一张空图。兵额、行军路
线、登陆点、战损和当地社会的反应，并非本纪短句所能逐项保存；地方群体更没
有借这条朝廷记录为自己留下完整声音。可以把行动放回海峡的风、潮、给水与补
给困难中理解，却不能由后来的海图倒推出吴军的每一步，或将一次军事行动写成
长期、无阻的控制。

珠崖、儋耳的记载同卫温远航、辽东使团构成的不是“孙吴海权”清单，而是几种
彼此不同的海上试探：有的是追逐传闻中的岛屿，有的是让礼物抵达政治盟点，有
的是向海峡另一岸派兵。它们所需的船、人、粮和消息都可能来自同一套沿江、沿
海的供应网络，却面对不同的风险和不同的对方。若只从建业的命令出发，便会看
不见港口劳动者、沿岸居民与岛屿社会并不是被动的地理；若只把失败归为风浪，
也会看不见朝廷如何因自身的战事、财政和人事不断改变远航的目的。

到二六三年，吴丞相兴建议取屯田万人为兵。这一条与蜀汉覆亡的年份相接：西方
战局的变化使吴廷重新估量可用于防务的人手。记载中的“万人”应当作为官方叙
事所列的动员数保留，不能据此推算所有屯田户或全境兵源；它仍然指出一个尖锐
的选择。把耕作系统中的人转为军人，可能在短期内增加可调的兵力，却会削弱原
先为粮食、运输与地方维持提供的人手。国家以征调缓解眼前的防务压力，同时把
下一季的供给和家庭安排变成新的问题。

这一选择使走马楼的近景重新有了位置。临湘木简并不记录二六三年这一次具体征
兵，也不能拿来证明命令如何在全吴执行；但它让我们理解为何“取人”绝不是从
一张静止名册上划掉一个数字。有人被抽离，田、租、船运、亲属责任和基层吏役
都须重新接合。孙吴在晚期仍能发出命令、修筑通道、集结军队，也正因如此更应
追问这些命令把不确定性留在何处。它的制度所得，是在江河和海峡之间维持一个
可调度的政权；它的制度代价，则是将防务、工程和财政反复压回地方社会那些最
难被史书完整记录的连接处。

\subsection{从一县的简到一国的限度}

把临湘、句容、濡须、会稽、辽东航路和珠崖、儋耳放在同一条阅读线上，并不是
说这些地方受同一套日常制度整齐支配。相反，它们之所以能够共同说明孙吴，正
在于国家每向外延伸一次，就必须换一种媒介：在临湘，它需要书手和可核的田租；
在山地，它需要置郡、派守和难以看清的协商或冲突；在江面，它需要船、渡口和
比对手更早抵达的信息；在外海，它依赖港口传闻、补给与对岸愿不愿接受联络。
这些媒介彼此牵连，却不能互相替代。能调船并不等于能使山地人入籍，能下诏减
赋也不等于能保证一座县廷立即改写所有欠负。

这也说明为什么孙吴的地方社会不能只写成“大族、山越与流民”的固定分类。文
书所见的是被行政辨认的一部分人，正史所见的是被朝廷认为值得记下的一部分行
动，考古所见又取决于保存和发掘的机会。三种视野都留下大量没有名字的人：在
仓门前等候交接者，在道路工程旁供给或被征者，在渡口提供消息者，在海峡另一
岸面对军船者。他们并不是可以由后世小说填补的空格。更严谨的写法是承认其不
可见，同时指出国家恰须经过他们的劳动、知识与土地关系才能把一次命令变成一
项结果。

孙吴最终在二八〇年被西晋的多路水陆进军所终结，不能倒过来使此前所有江防、
港口和文书都显得徒然。它们曾使一个以长江为轴的政权能够征收、转运、守备并
同远方接触；西晋的上游舟师和并进路线之所以构成威胁，也正因为它们冲击了这
些已被长期维持的节点。反过来，孙吴末期把屯田人转为兵的选择也不应被简单读
成“衰亡的证据”：它是在既有资源和新的威胁之间作出的短期处置，代价则落在
难以由史书逐项追踪的供给与家庭关系上。

读这一段江河史，最后不应寻找一幅孙吴社会的完成图。临湘的简牍、吴廷本纪和
城址、航路的研究各自保存的是不同距离上的痕迹。它们合在一起，能让我们看见
一条更具体的因果链：朝廷要维护江防和远近关系，便须取得粮、人、船与信息；
取得它们必须通过户籍、田租、市场、道路、港口和地方中介；每一次通过都会制
造拖欠、逃避、损耗、重新谈判或史料无法照亮的成本。国家并非浮在江面之上，
它也被这些不断需要重做的地方关系支撑，并最终受其限度约束。

\subsection{转运迟到时，问题落在谁的门前}

江河世界最容易被写成一张流动的地图：上游的木材顺流而下，稻米进入仓廪，舟
师在渡口集结，使者从港口离岸。真正困难的却是每一次流动都可能停在中途。一
场风雨能使船期错开，水位变化能使原先可用的泊位暂时不能卸货，修船所需的木
材和铁器也会同农具、房屋或其他工程争夺。国家文书通常在命令发出、工程完成
或军事行动的结果处停笔；县廷和港口面对的则是迟到以后怎样继续登记、怎样解
释损耗、怎样另找人手。没有足够材料，我们不能把这些情形编成某一年的完整危
机日志，但也不能因此把转运想成没有等待和折损的管道。

这种迟误会改变家庭与地方中介的位置。某一户的田租若尚未交清，未必只因其怠
惰，也可能同水路、收纳程序、劳役冲突或原有依附关系有关；不过简牍并不让后
人逐项判定原因。对基层吏员来说，模糊并不意味着可以不处理：责任须被暂时记
在某个户、田、仓或人名旁，直到下一次核验给出新的处置。对持有军令的人来说，
同一迟误可能只是“未集”的数字，却可能使沿岸社会失去一段可耕作、可交易或
照料家人的时间。正是在这个不对称里，财政、运输和战争不再是三条平行的叙事，
而成为彼此推挤的日常。

孙吴统治者有时用宽赋、省役或亲耕的语言回应这种压力，有时又以工程、铸钱和
征兵把资源再度集中。不要将前者读为毫无成本的仁政，也不要将后者化成单纯的
掠夺标签。每一种安排都试图解决眼前的一个瓶颈：使粮食可得、军队可动、市场
可通或边防可守；每一种安排也会把成本移给其他节点。对三国地方社会最有解释
力的，因而不是“江南已开发”或“江南尚蛮荒”的二分，而是追问在某一项具体
调度中，谁拥有改期、躲避、担保或重新协商的余地，谁只能接受被写入下一根简
的责任。

因此，租米、邸阁、海船和山地置郡也不宜各自退回“经济”“交通”“外交”或
“民族关系”的独立章节。它们在孙吴的时间里经常交会。一条可通的路会让市场
和守军同时靠近，也会带来新的劳役；一次海上使团需要珍货、船只和港口，同时
又把地方供给卷入远方政权的选择；一项户籍或田租手续原本属于县廷，却可能在
征兵、治水或转输时成为重新衡量家庭责任的根据。将这些关系写成同一过程，并
不是取消它们的差异，而是拒绝让宫廷年表把地方生活切碎成没有后果的背景。

材料的限度在这里不只是谨慎的尾注。它改变了我们能够提出的问题：我们能问一
次诏令为何需要经过县廷，不能据此宣布它在全境生效；我们能问一次航行为何会
暴露补给问题，不能为船队补出一条精确的现代航线；我们能问“山越”这一称谓
怎样服务于吴廷的行政和军事语言，不能让它替那些人群说出从未保存的共同意志。
保留这些界限，才使临湘的木简、吴主的本纪和海上、山地的零星记录彼此照亮，
而不是被拼装成一幅看似完整、实则越过证据承载力的孙吴社会图。
\label{reader:wu-river-society:end}
